\subsection{The Compositional Bottleneck in Diffusion Transformers}
Improvements in compositional accuracy are most pronounced for vertical relations---CogVLM raises accuracy on ``above'' from 72.73\% to 90.00\% and on ``below'' from 70.00\% to 90.00\% (Table~\ref{tab:accuracy})---providing a quantitative refutation of the prevailing assumption that generative models require full-model fine-tuning to improve structural adherence. Instead, our results suggest that the Diffusion Transformer (DiT) backbone in SD3.5 possesses a latent generative ``competence'' that remains underutilized due to impoverished conditioning signals from standard text encoders.

While standard CLIP-L is known to exhibit ``bag-of-words'' behavior that prioritizes mere object presence over relational syntax, the RedCaps CLIC encoder leverages Reddit’s alt-text, which is frequently authored for accessibility purposes to provide granular spatial grounding for the visually impaired. This transition from noisy, web-crawled captions to descriptive, accessibility-focused text allows the text encoder to move beyond simple object-token mapping and toward a truly relational embedding space.

The CogVLM CLIC encoder shows a complementary, but distinct, behavioral profile. In our evaluation, CogVLM delivers especially strong gains on challenging directional composition (notably ``to the right of''), indicating that encoder replacement can recover spatial distinctions that the baseline text encoder under-specifies. At the representation level, CogVLM exhibits consistently larger Frobenius norms than RedCaps across attention and MLP components, with the largest relative increase in \texttt{mlp.fc2.weight} and a clear elevation in Key, Value, and Output projections. Together, these trends suggest that CogVLM supplies a stronger, more targeted compositional conditioning signal to the same SD3.5 DiT backbone, bypassing the need for computationally expensive backbone retraining.

\subsection{The ``Reasoning-Output Paradox'' in VLMs}
Our evaluation of Qwen3-VL-30B revealed a fundamental discrepancy between a model's internal reasoning traces and its final classification output, a phenomenon we term the \textit{Reasoning-Output Paradox}.

\subsubsection{Internal Representation vs. Classification Head} 
In multiple instances, Qwen3-VL generated a logically sound and physically accurate reasoning trace, yet concluded with an incorrect spatial answer. This suggests that the model's internal spatial awareness is significantly more advanced than its final classification head's ability to map that awareness to a rigid, categorical label.

\subsubsection{Inverse Verbosity and Ambiguity}
The tendency for Qwen3-VL to label standard SD3.5 outputs as ``Ambiguous'' more frequently in Inverse Prompting mode highlights a secondary failure mode: when the generated image itself lacks coherent spatial stratification, the VLM defaults to over-verbose descriptions that ultimately fail to trigger a definitive match with the target prompt's logic.

\subsection{Deep Dive into Failure Modes: Semantic Priors vs. Structural Nouns}
The exploratory phase of this project uncovered generation failures that point to a fundamental semantic hierarchy problem within the diffusion process itself.

\subsubsection{The ``Pizza Flyer'' Phenomenon}
When the model generates a literal pizza instead of a paper flyer, it visually illustrates that high-probability object tokens (e.g., ``pizza'') frequently override low-probability structural nouns (e.g., ``flyer''). This suggests that the diffusion transformer's cross-attention layers remain heavily biased toward the most visually salient tokens in a prompt, regardless of their actual grammatical or modifying role.

\subsubsection{Left-Right Asymmetry}
The persistent confusion between ``left'' and ``right'', as well as ``above'' and ``below'', across all evaluated models, even after surgical encoder replacement, indicates a potential backbone-level bias. This demonstrates that even 70B+ parameter models struggle with directional absolute positioning, possibly due to the mirror-invariant properties of most underlying image-processing kernels and data augmentation strategies.

\subsection{Surgical Insight: Weight Divergence in Layers}
Our layer-wise analysis of the encoder replacement shows a broad increase in Frobenius norms from the early to late transformer blocks (layers 0--11), with the largest magnitudes concentrated in the final layers. This specific pattern indicates that the encoder substitution most meaningfully shifts the higher-level text representations utilized for DiT conditioning. Component-wise, CogVLM exhibits higher norms than RedCaps across both attention and MLP weights, with the largest relative increase located in \texttt{mlp.fc2.weight}. Within the attention mechanism itself, Key, Value, and Output projections are all elevated. These results are highly consistent with a modified, finer-grained token-interaction profile introduced after the encoder replacement.

\subsection{Interpreting the GenEval2 TIFA Scores}
The GenEval2 Soft-TIFA results reinforce, and in one case complicate, the picture drawn by our spatial-relation benchmarks. The 14 pp gap between the paper-reported SD3.5-Large score (69.7\%) and our Turbo baseline (55.5\%) is itself a finding: it quantifies how much compositional accuracy is sacrificed by the Turbo distillation, and suggests that our modular encoder intervention is applied to an already-constrained generative capacity.

Within our Turbo variants, both CLIC variants improve over the baseline on AM and GM, confirming that compositionally-aware text encoders provide a reliable, if incremental, grounding benefit across a diverse 800-prompt benchmark---not merely on the curated relational prompts we designed. The difference between CogVLM leading on AM while RedCaps leads on GM suggests the two variants differ in character: CogVLM more reliably handles individual visual atoms, while RedCaps is more consistent across entire prompts. One plausible explanation is that CogVLM's richer, higher-norm attention representations improve fine-grained token grounding but occasionally fail whole prompts, whereas RedCaps offers a more uniform grounding signal. This finding motivates future exploration of ensemble or skill-specific encoder selection strategies.

\subsection{Limitations and Future Outlook}
Despite the demonstrable success of our modular intervention, certain constraints remain:

\begin{enumerate}
    \item \textbf{Logical Negation:} Despite some measurable improvement with the concept of ``without'' when using the RedCaps and CogVLM SD3.5 variants, the broader problem of strict logical negation, previously documented in earlier Stable Diffusion models, remains largely unsolved.
    \item \textbf{Vague Phrasing:} Relational phrases such as ``far from'', ``closeby'', and ``on'' proved to be highly subjective during evaluation, leading to the lowest levels of agreement between the automated VLM predictions and our manual human annotations.
    \item \textbf{Hardware Efficiency:} While we optimized the evaluation pipeline for NVIDIA A100 clusters, the sheer computational cost of large-scale VLM assessment remains a significant overhead bottleneck for rapid, iterative testing in generative model development.
\end{enumerate}

\subsection{Conclusion}
This research establishes that compositionality is not merely an emergent property of model scale, but rather a direct consequence of training data quality and text-encoder grounding. By replacing the standard CLIP-L encoder with compositionally-aware variants like RedCaps and CogVLM, we achieved measurable improvements in spatial reasoning on select relational phrases, a finding corroborated by the broader GenEval2 Soft-TIFA evaluation. Ultimately, this ``weight surgery'' approach provides a robust, production-ready framework for enhancing the spatial reliability of text-to-image models without incurring the prohibitive computational costs associated with full-model retraining.